\SGAchapter{VIII}{双扩张的补充；群概形的双扩张的一般性质}{Compléments sur les biextensions. Propriétés générales des biextensions des schémas en groupes}
\begin{center}A. Grothendieck\end{center}
\SourcePage{227}
\SGAsection{0}{引言}
\textbf{0.1.} 我们在这里继续上一讲开始的\Term{biextension}{双扩张}的\Term{generic}{一般}研究。首先，在\Term{generic}{一般}基\Term{topos}{拓扑斯}的情形给出一些补充，特别是详尽展开与\Term{biextension}{双扩张}相联系的\Term{pairing}{配对}的\Term{form}{形式}理论（第2节）；这些\Term{pairing}{配对}在\Term{abelian-scheme}{阿贝尔概形}理论中的重要性是众所周知的。另一方面，我们开始研究\Term{generic}{一般}\Term{group-scheme}{群概形}的\Term{biextension}{双扩张}；对我们来说，最有意义的情形是以\Term{multiplicative-group}{乘法群} $\mathbf G_m$ 为\Term{kernel}{核}的\Term{biextension}{双扩张}。不过，作为技术上的中间步骤，我们还不得不研究以其他类型的\Term{group-scheme}{群概形}为\Term{kernel}{核}的\Term{biextension}{双扩张}（第5、6节）。
\SourcePage{228}
对后文最有意义的结果是7.1：设 $P,Q$ 在一个\Term{generic-point}{泛点}为 $\eta$ 的\Term{trait}{离散赋值环谱}上\Term{smooth}{光滑}，该结果给出一些条件，保证 $(P_\eta,Q_\eta)$ 的一个以 $\mathbf G_{m,\eta}$ 为\Term{kernel}{核}的\Term{biextension}{双扩张}可以延拓为 $(P,Q)$ 的一个以 $\mathbf G_m$ 为\Term{kernel}{核}的\Term{biextension}{双扩张}，且在唯一\Term{isomorphism}{同构}意义下确定。这个结果在下一讲研究\Term{abelian-variety}{阿贝尔簇}的 \Term{neron-model}{Néron 模型}时会很方便。很可能不必搬出第VII、VIII讲这两块沉重的镇纸（本研讨班后文不再使用它们），也能以较小的代价完成那里的工作；不过，我们认为这两讲所展开的内容很可能还会在别处有用。

\textbf{0.2.} 当我们研究\Term{scheme}{概形} $S$ 上的\Term{group-scheme}{群概形}时，约定在 $S$ 的 fppf \Term{topos}{拓扑斯}中工作，即：取 $S$ 上\Term{locally-finitely-presented}{局部有限表示}且属于 $\mathcal U$ 的\Term{scheme}{概形}所成的 $\mathcal U$-\Term{site}{景}，赋予 fppf \Term{topology}{拓扑}，再取其关联\Term{topos}{拓扑斯}（SGA 3，第IV讲6.3）。读者也很容易\Term{kernel}{核}实，我们给出的大部分命题在 $S$ 上其他一些\Term{topology}{拓扑}中同样成立，而且与 fppf 版本本质\Term{equivalence}{等价}；这里要求\Term{site}{景}的底\Term{category}{范畴}至少包含 $S$ 上所有\Term{locally-finitely-presented}{局部有限表示}的\Term{scheme}{概形}。例如可以取 Zariski \Term{topology}{拓扑}、\Term{etale}{平展}\Term{topology}{拓扑}或 fpqc \Term{topology}{拓扑}。原因之一是：由\Term{descent}{下降}理论，$S$ 上\Term{locally-finitely-presented}{局部有限表示}的\Term{group-scheme}{群概形} $P$ 以\Term{group-scheme}{群概形} $\mathbf G_m$ 为\Term{kernel}{核}的\Term{group-extension}{扩张}所成的\Term{category}{范畴}，在\Term{equivalence}{等价}意义下不依赖所选\Term{topology}{拓扑}；\Term{biextension}{双扩张}也有同样的性质。宁取 fppf \Term{topology}{拓扑}（或 fpqc 等更细的\Term{topology}{拓扑}）而不取 Zariski \Term{topology}{拓扑}、\Term{etale}{平展}\Term{topology}{拓扑}等较粗\Term{topology}{拓扑}，其技术原因是：对每个整数 $n>0$，$\mathbf G_m$ 的自\Term{homomorphism}{同态} $n\operatorname{id}$ 都是 fppf \Term{epimorphism}{满态射}（因为它\Term{faithfully-flat}{忠实平坦}且\Term{finite-presentation}{有限表示}），但\Term{generic}{一般}不是 (Zar) 或 (ét) \Term{epimorphism}{满态射}。（注意：它是 (ét) \Term{epimorphism}{满态射}当且仅当 $n$ 与 $S$ 的各\Term{residue-characteristic}{剩余特征}互素；它是 (Zar) \Term{epimorphism}{满态射}当且仅当 $S$ 为空。）
\SourcePage{229}
\textbf{0.3.} 最后指出，大部分关于以 $\mathbf G_m$ 为\Term{kernel}{核}的\Term{group-extension}{扩张}或\Term{biextension}{双扩张}的命题，把这个\Term{group}{群}换成\Term{torus}{环面} $G$ 后仍成立。因为这样的 $G$ \Term{local-scheme}{局部}地（对于 fppf，甚至对于 (ét)）\Term{isomorphism}{同构}于 $\mathbf G_m^r$，所以这种推广本质上是平凡的。

\SGAsection{1}{任意拓扑斯上的若干特殊情形}
\textbf{1.1.} 在本节和下一节中，与第VII讲一样，$T$ 表示一个\Term{topos}{拓扑斯}，$P,Q,G$ 表示 $T$ 的三个\Term{abelian-group}{阿贝尔群}。因此有\Term{canonical-isomorphism}{典范同构}（第VII讲3.6.5）
\[\operatorname{Biext}^1(P,Q;G)\simeq\operatorname{Ext}^1(P\mathbin{\overset{\mathbf L}{\otimes}}Q;G).\tag{1.1.1}\]
回顾还有称为“Cartan \Term{isomorphism}{同构}”的\Term{canonical-isomorphism}{典范同构}
\[\begin{split}
\operatorname{Ext}^1(P\mathbin{\overset{\mathbf L}{\otimes}}Q,G)
&\simeq\operatorname{Ext}^1(P,\mathbf R\mathcal Hom(Q,G))\\
&\simeq\operatorname{Ext}^1(Q,\mathbf R\mathcal Hom(P,G)).
\end{split}\tag{1.1.2}\]
所有这些\Term{isomorphism}{同构}对三个变元 $P,Q,G$ 都是\Term{functoriality}{函子性}的。
由（1.1.1），考虑到 $P\mathbin{\overset{\mathbf L}{\otimes}}Q$ 的\Term{homology-object}{同调对象}为 $\mathcal Tor_i(P,Q)$，且它们在 $i\geq1$ 时为零，得到以下无限\Term{exact-sequence}{正合列}：
\[\begin{split}
0\longrightarrow\operatorname{Ext}^1(P\otimes Q,G)&\longrightarrow\operatorname{Biext}^1(P,Q;G)
\longrightarrow\operatorname{Hom}(\mathcal Tor_1(P,Q),G)\\
&\longrightarrow\operatorname{Ext}^2(P\otimes Q,G)
\longrightarrow\operatorname{Ext}^2(P\mathbin{\overset{\mathbf L}{\otimes}}Q,G)\longrightarrow\cdots .
\end{split}\tag{1.1.3}\]
另一方面，$\mathbf R\mathcal Hom(Q,G)$ 的\Term{cohomology-object}{上同调对象}为 $\mathcal Ext^i(Q,G)$，故（1.1.2）的第一个等式给出五项\Term{exact-sequence}{正合列}
\[\begin{split}
0\longrightarrow\operatorname{Ext}^1(P,\mathcal Hom(Q,G))&\longrightarrow\operatorname{Biext}^1(P,Q;G)
\longrightarrow\operatorname{Hom}(P,\mathcal Ext^1(Q,G))\\
&\longrightarrow\operatorname{Ext}^2(P,\mathcal Hom(Q,P))
\longrightarrow\operatorname{Ext}^2(P,\mathbf R\mathcal Hom(Q,G)),
\end{split}\tag{1.1.4}\]
\SourcePage{230}
它关联于通常的\Term{spectral-sequence}{谱序列}
\[\operatorname{Ext}^*(P,\mathbf R\mathcal Hom(Q,G))\Longleftarrow
\operatorname{Ext}^p(P,\mathcal Ext^q(Q,G)).\]
如果读者觉得需要，可以用自然的\Term{fully-faithful-functor}{充分忠实函子}
\[\operatorname{EXT}(P\otimes Q,G)\longrightarrow\operatorname{BIEXT}(P,Q;G)\tag{1.1.5}\]
（相应地用
\[\operatorname{EXT}(P,\mathcal Hom(Q,G))\longrightarrow\operatorname{BIEXT}(P,Q;G)\tag{1.1.6}\]
）明确描述（1.1.3）（相应地（1.1.4））的第一个\Term{arrow}{箭头}；这里 $\operatorname{EXT}$ 和 $\operatorname{BIEXT}$ 分别表示不难猜出的\Term{group-extension}{扩张}\Term{category}{范畴}和\Term{biextension}{双扩张}\Term{category}{范畴}。
当然，\Term{isomorphism}{同构} $\operatorname{Biext}^1(P,Q;G)\simeq\operatorname{Ext}^1(Q,\mathbf R\mathcal Hom(P,G))$ 给出与（1.1.3）类似但交换了 $P,Q$ 地位的\Term{exact-sequence}{正合列}；它的第一个\Term{arrow}{箭头}可以用与（1.1.6）类似的\Term{fully-faithful-functor}{充分忠实函子}
\[\operatorname{EXT}(Q,\mathcal Hom(P,G))\longrightarrow\operatorname{BIEXT}(P,Q;G)\tag{1.1.7}\]
来明确描述。

\textbf{1.2.} 假设
\[\mathcal Tor_1(P,Q)=0,\tag{1.2.1}\]
则（1.1.3）给出\Term{canonical-isomorphism}{典范同构}
\[\operatorname{Biext}^1(P,Q;G)\simeq\operatorname{Ext}^1(P\otimes Q,G).\tag{1.2.2}\]
例如，$P$ 或 $Q$ 是\Term{flat}{平坦} $\mathbf Z$-\Term{module}{模}时就是如此。显然，可以更精确地说：在条件（1.2.1）下，\Term{functor}{函子}
\SourcePage{231}
（1.1.5）是\Term{equivalence-of-categories}{范畴等价}。

\textbf{1.3.} 假设
\[P\otimes Q=0,\tag{1.3.1}\]
则（1.1.3）给出\Term{canonical-isomorphism}{典范同构}
\[\operatorname{Biext}^1(P,Q;G)\simeq\operatorname{Hom}(\mathcal Tor_1(P,Q),G).\tag{1.3.2}\]
对我们最有意义的满足（1.3.1）的情形是：存在整数 $n>0$，使得
\[nP=P,\qquad Q=\varinjlim_i {}_{n^i}Q,\tag{1.3.3}\]
即 $P$ 是 $n$-\Term{divisible}{可除}的，而 $Q$ 是被 $n$ 的\Term{nilpotent}{幂零}化的\Term{subgroup}{子群}的\Term{limit}{极限}。照常记
\[T_n(P)=({}_{n^i}P)_{i\geq0}\quad\text{（投射系）},\qquad
{}_{n^\infty}P=\varinjlim_i{}_{n^i}P,\]
则由（1.3.3）可得比（1.3.1）更精确的公式
\[\begin{split}
P\otimes Q&=0,\\
\mathcal Tor_1(P,Q)&\simeq\varinjlim_i({}_{n^i}P\otimes{}_{n^i}Q)
\simeq T_n(P)\otimes Q\simeq{}_{n^\infty}P\otimes T_n(Q),
\end{split}\tag{1.3.4}\]
因而\Term{isomorphism}{同构}（1.3.2）在这里具有更具体的\Term{form}{形式}：
\[\begin{split}
\operatorname{Biext}^1(P,Q;G)&\simeq\operatorname{Hom}(T_n(P),D(Q))\\
&\simeq\operatorname{Hom}(Q,D(T_n(P)))\\
&\simeq\operatorname{Hom}(T_n(P)\otimes T_n(Q),T_n(G)).
\end{split}\tag{1.3.5}\]
下面说明这些公式的意义。

 a）在（1.3.4）出现的\Term{inductive-system}{归纳系}中，\Term{transition-homomorphism}{过渡同态}
\[{}_{n^i}P\otimes{}_{n^i}Q\longrightarrow{}_{n^j}P\otimes{}_{n^j}Q\qquad(i\leq j)\]
\SourcePage{232}
可如下描述：先借助 $n^{j-i}\otimes\operatorname{id}$ 把第一个成员等同于 ${}_{n^j}P\otimes{}_{n^i}Q$，再使用\Term{homomorphism}{同态}
\[\operatorname{id}\otimes\operatorname{inj}:{}_{n^j}P\otimes{}_{n^i}Q\longrightarrow{}_{n^j}P\otimes{}_{n^j}Q,\]
其中 $\operatorname{inj}:{}_{n^i}Q\longrightarrow{}_{n^j}Q$ 表示\Term{inclusion}{包含映射}。

 b）定义
\[T_n(P)\otimes Q\overset{\mathrm{def}}=\varinjlim_i T_n(P)\otimes{}_{n^i}Q,\]
其中
\[T_n(P)\otimes M\overset{\mathrm{def}}={}_{n^i}P\otimes M\quad\text{若 }n^iM=0.\]
若再假设 $nQ=Q$（从而 $T_n(Q)$ 与 $T_n(P)$ 一样是严格\Term{projective-system}{投射系}），便可类似地说明（1.3.4）所写第三项的意义。（1.3.4）的最后一项是\Term{group}{群}的\Term{projective-system}{投射系}之间的一个 $\operatorname{Hom}$。

 c）字母 $D$ 表示 $\mathcal Hom(-,G)$。\Term{object}{对象} $D(Q)$ 定义为\Term{group}{群}的\Term{projective-system}{投射系}
\[D(Q)=(D({}_{n^i}Q))_{i\geq0}=(\mathcal Hom({}_{n^i}Q,G))_{i\geq0};\]
（1.3.5）的第二个成员是\Term{group}{群}的\Term{projective-system}{投射系}的\Term{homomorphism}{同态}\Term{group}{群}。最后一个成员的解释更简单，只需置
\[D(T_n(P))=\varinjlim_i D({}_{n^i}P),\]
这个\Term{inductive-system}{归纳系}的\Term{transition-morphism}{过渡态射}当然由 $T_n(P)$ 的\Term{transition-morphism}{过渡态射}
\SourcePage{233}
$n^{j-i}:{}_{n^j}P\longrightarrow{}_{n^i}P$ \Term{transpose}{转置}而来。
我们把（1.3.4）所述 $\mathcal Tor_1(P,Q)$ 的计算及其推论（1.3.5）的细节留给读者。只需注意，我们用到由（1.3.3）的第二个假设\Term{derive}{导出}的关系
\[\mathcal Tor_p(P,Q)\simeq\varinjlim_i\mathcal Tor_p(P,{}_{n^i}Q);\]
当 $p=0$ 时立即得到 $P\otimes Q=0$；这也把 $\mathcal Tor_1(P,Q)$ 的计算归结为各 $\mathcal Tor_1(P,{}_{n^i}Q)$ 及其\Term{transition-morphism}{过渡态射}的计算。因此，当 $Q$ 被一个整数 $m=n^i$ \Term{annihilate}{零化}且 $mP=P$ 时，只须明确给出一个\Term{isomorphism}{同构}
\[\mathcal Tor_1(P,Q)\xleftarrow{\sim}{}_{m}P\otimes{}_{m}Q
\qquad(mP=P,\ mQ=0).\tag{1.3.6}\]
这有多种做法，其中一种是使用与\Term{exact-sequence}{正合列}
\[0\longrightarrow{}_{m}P\longrightarrow P\xrightarrow{m}P\longrightarrow0\]
关联的 $\mathcal Tor$ \Term{exact-sequence}{正合列}。也可以使用下文的\Term{homomorphism}{同态}（2.1.1）；此时\Term{transition-morphism}{过渡态射}正是 a）所述的\Term{morphism}{态射}这一事实，就是2.1.11。无论采用哪一种方法，（1.3.6）的选择都归结为一个\Term{signature}{符号}的选择，参见2.1.10。

\textbf{1.4.} 假设
\[\mathcal Hom(Q,G)=0,\tag{1.4.1}\]
则（1.1.4）给出\Term{canonical-isomorphism}{典范同构}
\[\operatorname{Biext}^1(P,Q;G)\xrightarrow{\sim}\operatorname{Hom}(P,\mathcal Ext^1(Q,G)).\tag{1.4.2}\]
\SourcePage{234}
特别地，当 $Q$ 是\Term{abelian-scheme}{阿贝尔概形}且 $G=\mathbf G_m$ 时就是如此，参见第VII讲1.3.8。因此，关于 $P$ 的\Term{functor}{函子} $P\mapsto\operatorname{Biext}^1(P,Q;G)$ 由 $Q^*=\mathcal Ext^1(Q,G)$ 表示；特别得到 $(Q^*,Q)$ 的一个以 $G$ 为\Term{kernel}{核}的\Term{canonical}{典范}\Term{biextension}{双扩张} $E^Q$，在唯一\Term{isomorphism}{同构}意义下确定（因为 $E^Q$ 的\Term{automorphism-group}{自同构群}为 $\operatorname{Hom}(P,\mathcal Hom(Q,G))=0$）。我们可以把这个\Term{biextension}{双扩张}，或 $(Q,Q^*)$ 的\Term{symmetric}{对称}\Term{biextension}{双扩张} $sE^Q$，称为 $(Q^*,Q)$ 或 $(Q,Q^*)$ 的以 $G$ 为\Term{kernel}{核}的 Weil \Term{biextension}{双扩张}。

\textbf{1.5.} 假设
\[\operatorname{Hom}(P,\mathcal Ext^1(Q,G))=0.\tag{1.5.1}\]
特别地，若 $\mathcal Ext^1(Q,G)=0$，这一条件成立（例如 $Q$ 是\Term{finite-locally-free}{有限局部自由}\Term{group-scheme}{群概形}或\Term{multiplicative-type-group}{乘法型群}（SGA 3，第IX讲），且 $G=\mathbf G_m$；参见下文3.4）。于是（1.1.4）给出\Term{canonical-isomorphism}{典范同构}
\[\operatorname{Biext}^1(P,Q;G)\xleftarrow{\sim}\operatorname{Ext}^1(P,\mathcal Hom(Q,G)),\tag{1.5.2}\]
这里 $\mathcal Hom(Q,G)$ 起到 $Q$ 的一个“\Term{dual-group}{对偶群}”（相对于 $G$）的\Term{action}{作用}。更精确地，由（1.5.2）可知\Term{fully-faithful-functor}{充分忠实函子}（1.1.6）是\Term{equivalence-of-categories}{范畴等价}。

\textbf{1.6.} 假设同时有
\[\mathcal Hom(P,G)=0,\qquad\operatorname{Hom}(P,\mathcal Ext^1(Q,G))=0.\tag{1.6.1}\]
根据1.4，这两个关系分别给出\Term{isomorphism}{同构}
\[\operatorname{Biext}^1(P,Q;G)\xrightarrow{\sim}\operatorname{Hom}(Q,\mathcal Ext^1(P,G)),\]
\[\operatorname{Biext}^1(P,Q;G)\xleftarrow{\sim}\operatorname{Ext}^1(P,\mathcal Hom(Q,G));\]
\SourcePage{235}
把它们复合，得到\Term{canonical-isomorphism}{典范同构}
\[\operatorname{Ext}^1(P,\mathcal Hom(Q,G))\xrightarrow{\sim}
\operatorname{Hom}(Q,\mathcal Ext^1(P,G))\quad(\simeq\operatorname{Biext}^1(P,G;G)),\tag{1.6.2}\]
它起到\Term{duality}{对偶}公式的\Term{action}{作用}。还应注意，（1.6.1）的第一个关系显然蕴含 $\operatorname{Hom}(P\otimes Q,G)=0$，所以 $P$ 以 $\mathcal Hom(Q,G)$ 为\Term{kernel}{核}的（交换）\Term{group-extension}{扩张}的\Term{category}{范畴} $\operatorname{EXT}(P,\mathcal Hom(Q,G))$ 是\Term{rigid-category}{刚性}的（\Term{object}{对象}的每个\Term{automorphism}{自同构}都是恒等）；$\operatorname{BIEXT}(P,Q;G)$ 也一样。因此，只要知道\Term{object}{对象}的\Term{isomorphism-class}{同构类}\Term{group}{群}，就知道这些\Term{category}{范畴}的结构；而根据（1.6.2），这个\Term{group}{群}\Term{canonical-isomorphism}{典范同构}于 $\operatorname{Hom}(Q,\mathcal Ext^1(P,Q))$。
还应注意，无论是否假设（1.6.1），映射（1.6.2）总有定义，而且可以不经过 $\operatorname{Biext}^1(P,Q;G)$，直接用自然\Term{pairing}{配对}
\[Q\times Q'\longrightarrow G\qquad(Q'=\mathcal Hom(Q,G))\]
来描述：把它解释为\Term{homomorphism}{同态} $Q\longrightarrow\mathcal Hom(Q',G)$，便得到复合
\[\operatorname{Ext}^1(P,Q')\times Q\longrightarrow
\operatorname{Ext}^1(P,Q')\times\mathcal Hom(Q',G)\longrightarrow\mathcal Ext^1(P,G),\tag{1.6.3}\]
这就重新得到（1.6.2）。关于1.6适用的一个有用情形，参见下文3.7。

\SGAsection{2}{由双扩张定义的配对}
\textbf{2.1.} 设 $P,Q$ 是 $T$ 上的两个\Term{abelian-group}{阿贝尔群}，$n>0$ 为整数。回顾，${}_nP$ 表示乘以 $n$ 的\Term{homomorphism}{同态} $n\operatorname{id}_P$ 的\Term{kernel}{核}。我们将定义一个\Term{canonical}{典范}\Term{homomorphism}{同态}
\SourcePage{236}
\[\alpha=\alpha^n_{P,Q}:{}_nP\otimes{}_nQ\longrightarrow\mathcal Tor_1(P,Q),\tag{2.1.1}\]
它是复合
\[{}_nP\otimes{}_nQ\longrightarrow\mathcal Tor_1({}_nP,{}_nQ)\longrightarrow\mathcal Tor_1(P,Q),\tag{2.1.2}\]
其中第一个\Term{arrow}{箭头}为 $\alpha^n_{{}_nP,{}_nQ}$，第二个\Term{arrow}{箭头}由\Term{inclusion}{包含映射} ${}_nP\longrightarrow P$、${}_nQ\longrightarrow Q$ \Term{derive}{导出}。因此，只须在 $nP=0,nQ=0$，即 $P,Q$ 是 $\mathbf Z/n\mathbf Z$-\Term{module}{模}时定义（2.1.1）。先只假设 $nQ=0$。取 $P$ 作为 $\mathbf Z$-\Term{module}{模}的\Term{flat-resolution}{平坦解消} $L_\bullet$，得到\Term{canonical-isomorphism}{典范同构}
\[P\mathbin{\overset{\mathbf L}{\otimes}}Q\simeq L_\bullet\otimes_{\mathbf Z}Q
\simeq(L_\bullet\otimes_{\mathbf Z}\mathbf Z/n\mathbf Z)\otimes_{\mathbf Z/n\mathbf Z}Q;\]
也可将它写为\Term{derived-category}{导出范畴}中的\Term{isomorphism}{同构}
\[P\mathbin{\overset{\mathbf L}{\otimes}}Q\simeq
(P\mathbin{\overset{\mathbf L}{\otimes}_{\mathbf Z}}\mathbf Z/n\mathbf Z)
\mathbin{\overset{\mathbf L}{\otimes}_{\mathbf Z/n\mathbf Z}}Q.\tag{2.1.3}\]
它给出一个\Term{homological}{同调型}的“Künneth”\Term{spectral-sequence}{谱序列}（参见 SGA 1，第IV讲5.2）：
\[\mathcal Tor_*^{\mathbf Z}(P,Q)\Longleftarrow
E^2_{p,q}=\mathcal Tor_p^{\mathbf Z/n\mathbf Z}(\mathcal Tor_q^{\mathbf Z}(P,\mathbf Z/n\mathbf Z),Q).\tag{2.1.4}\]
回顾有\Term{canonical-isomorphism}{典范同构}
\[\mathcal Tor_0^{\mathbf Z}(P,\mathbf Z/n\mathbf Z)=P_n\overset{\mathrm{def}}=P/nP,
\qquad\mathcal Tor_1^{\mathbf Z}(P,\mathbf Z/n\mathbf Z)\simeq{}_nP,\]
\[\mathcal Tor_i^{\mathbf Z}(P,\mathbf Z/n\mathbf Z)=0\quad(i\geq2),\]
这是用 $(\mathbf Z/n\mathbf Z)_T$ 的如下\Term{canonical}{典范}\Term{flat-resolution}{平坦解消}计算这些 $\mathcal Tor$ 得到的：
\[0\longrightarrow\underbrace{\mathbf Z_T}_{\mathbf L_1}\xrightarrow{n}
\underbrace{\mathbf Z_T}_{\mathbf L_0}\longrightarrow0.\]
\SourcePage{237}
\Term{augmentation}{增广} $\mathbf L_0\longrightarrow(\mathbf Z/n\mathbf Z)_T$ 把 $1$ 映到 $1$。因此，\Term{spectral-sequence}{谱序列}（2.1.4）本质上归结为\Term{exact-sequence}{正合列}
\[\begin{split}
0\longrightarrow\mathcal Tor_2^{\mathbf Z/n\mathbf Z}(P_n,Q)&\longrightarrow{}_nP\otimes Q
\xrightarrow{\alpha}\mathcal Tor_1^{\mathbf Z}(P,Q)\\
&\longrightarrow\mathcal Tor_1^{\mathbf Z/n\mathbf Z}(P_n,Q)\longrightarrow0,
\end{split}\tag{2.1.5}\]
其中 $P_n=P\otimes_{\mathbf Z}\mathbf Z/n\mathbf Z=P/nP$。所求\Term{homomorphism}{同态} $\alpha$ 就是这个\Term{exact-sequence}{正合列}中的\Term{edge-homomorphism}{边缘同态} $E^2_{0,1}\longrightarrow\mathcal Tor_1^{\mathbf Z}$。该\Term{exact-sequence}{正合列}显然对 $P,Q$ 都是\Term{functoriality}{函子性}的；对 $P$ 的\Term{functoriality}{函子性}表明，所考虑的\Term{homomorphism}{同态}确实是复合
\[{}_nP\otimes Q\longrightarrow\mathcal Tor_1({}_nP,Q)\longrightarrow\mathcal Tor_1(P,Q),\]
第一个\Term{homomorphism}{同态}是按上述方法、将 $(P,Q)$ 换成 $({}_nP,Q)$ 而定义的 $\alpha$，第二个\Term{homomorphism}{同态}来自\Term{inclusion}{包含映射} ${}_nP\longrightarrow Q$。

\textbf{2.1.6.} 因此，对任意 $P,Q$，这确实定义了满足\Term{resolution}{分解}条件（2.1.2）的\Term{homomorphism}{同态}（2.1.1）。此外，若 $nQ=0$，则\Term{exact-sequence}{正合列}（2.1.5）表明：只要 $P_n$ 在 $\mathbf Z/n\mathbf Z$ 上\Term{flat}{平坦}（特别地，若 $P_n=0$，即 $nP=P$，即 $P$ 是 $n$-\Term{divisible}{可除}的，或若 $P$ 本身是\Term{flat}{平坦} $\mathbf Z/n\mathbf Z$-\Term{module}{模}），或者 $Q$ 是\Term{flat}{平坦} $\mathbf Z/n\mathbf Z$-\Term{module}{模}，$\alpha$ 就是\Term{isomorphism}{同构}。

\textbf{2.1.7.} 具体计算（2.1.5）中的 $\alpha$。只要对 $Q$ 的每个\Term{section}{截面} $q$（定义在 $T$ 的\Term{object}{对象} $S$ 上；为简单起见，设 $S$ 是\Term{final-object}{终对象}），知道相应的\Term{homomorphism}{同态} ${}_nP\longrightarrow\mathcal Tor_1^{\mathbf Z}(P,Q)$，就知道 $\alpha$。而 $q$ 可以看成 $(\mathbf Z/n\mathbf Z)_T$ 的\Term{section}{截面} $1$ 在唯一确定的\Term{homomorphism}{同态} $u_q:(\mathbf Z/n\mathbf Z)_T\longrightarrow Q$ 下的像；于是 $\alpha$ 对 $Q$ 的\Term{functoriality}{函子性}
\SourcePage{238}
把我们归结到 $Q$ 是\Term{constant-sheaf}{常值层} $(\mathbf Z/n\mathbf Z)_T$ 且 $q$ 是它的\Term{unit-section}{单位截面}的情形。这时，按（2.1.4）之后的说明将 $\mathcal Tor_1^{\mathbf Z}(P,Q)=\mathcal Tor_1^{\mathbf Z}(P,\mathbf Z/n\mathbf Z)$ 等同于 ${}_nP$，所求\Term{homomorphism}{同态}就是恒等。换句话说，\Term{edge-homomorphism}{边缘同态}
\[E^2_{0,q}=\mathcal Tor_q^{\mathbf Z}(P,\mathbf Z/n\mathbf Z)
\otimes_{\mathbf Z/n\mathbf Z}\mathbf Z/n\mathbf Z
\longrightarrow\mathcal Tor_q^{\mathbf Z}(P,\mathbf Z/n\mathbf Z)\]
是恒等。这直接由定义得到（把 $\mathbf Z$ 及其\Term{quotient}{商} $\mathbf Z/n\mathbf Z$ 换成 $T$ 上任意\Term{ring}{环} $A$ 及其任意\Term{quotient-ring}{商环} $A/I$，类似的断言仍成立）。

\textbf{2.1.8.} 也可对 ${}_nP$ 的每个\Term{section}{截面} $p$ 和 ${}_nQ$ 的每个\Term{section}{截面} $q$（定义在 $T$ 的\Term{object}{对象} $S$ 上，为简化记号，仍取 $S$ 为\Term{final-object}{终对象}），描述 $\mathcal Tor_1^{\mathbf Z}(P,Q)$ 的\Term{section}{截面} $\alpha(p,q)$，以确定（2.1.1）的\Term{homomorphism}{同态} $\alpha$。如前，利用 $\alpha$ 对 $P,Q$ 的\Term{functoriality}{函子性}，可归结到
\[P=Q=(\mathbf Z/n\mathbf Z)_T,\]
$p,q$ 都是 $(\mathbf Z/n\mathbf Z)_T$ 的\Term{unit-section}{单位截面}的情形，也就是确定\Term{homomorphism}{同态}
\[\alpha:(\mathbf Z/n\mathbf Z)_T\otimes(\mathbf Z/n\mathbf Z)_T
\longrightarrow\mathcal Tor_1^{\mathbf Z}((\mathbf Z/n\mathbf Z)_T,(\mathbf Z/n\mathbf Z)_T).\]
2.1.7的描述表明，它就是\Term{isomorphism}{同构}
\[(\mathbf Z/n\mathbf Z)_T\xrightarrow{\sim}
\mathcal Tor_1^{\mathbf Z}((\mathbf Z/n\mathbf Z)_T,(\mathbf Z/n\mathbf Z)_T),\tag{2.1.8.1}\]
由 $\mathcal Tor_1$ 中第一个因子 $P=(\mathbf Z/n\mathbf Z)_T$ 的\Term{canonical}{典范}\Term{flat-resolution}{平坦解消}
\[0\longrightarrow\mathbf Z_T\xrightarrow{n}\mathbf Z_T\longrightarrow0\tag{2.1.8.2}\]
\SourcePage{239}
\Term{derive}{导出}。

\textbf{2.1.9.} （2.1.1）的定义并不\Term{symmetric}{对称}；交换上述构造中 $P,Q$ 的地位，先归结到 $nP=0$，再使用 $Q$ 的\Term{flat-resolution}{平坦解消}，可以定义类似的\Term{homomorphism}{同态}
\[\beta^n_{P,Q}:{}_nP\otimes{}_nQ\longrightarrow\mathcal Tor_1(P,Q).\tag{2.1.9.1}\]
\Term{equivalence}{等价}地，也可以由下图的\Term{commutativity}{交换性}，利用 $\alpha$ 来定义 $\beta$：
\[\begin{tikzcd}
{}_nP\otimes{}_nQ\arrow[r,"\sim"]\arrow[d,"{\beta^n_{P,Q}}"']&{}_nQ\otimes{}_nP\arrow[d,"{\alpha^n_{Q,P}}"]\\
\mathcal Tor_1(P,Q)\arrow[r,"\sim"]&\mathcal Tor_1(Q,P)
\end{tikzcd}\tag{2.1.9.2}\]
其中水平\Term{arrow}{箭头}是\Term{symmetry-isomorphism}{对称同构}。于是 $\alpha,\beta$ 的关系为
\[\beta^n_{P,Q}=-\alpha^n_{P,Q}.\tag{2.1.9.3}\]
换句话说：
\begin{SGAresult}{\Term{lemma}{引理}}{2.1.10}
设 $P,Q$ 为两个\Term{commutative-group}{交换群}，$n>0$ 为整数。则下图\Term{anticommutativity}{反交换}，其中 $\alpha$ 定义于（2.1.1），水平\Term{arrow}{箭头}是\Term{symmetry-isomorphism}{对称同构}：
\[\begin{tikzcd}
{}_nP\otimes{}_nQ\arrow[r,"\sim"]\arrow[d,"{\alpha^n_{P,Q}}"']&{}_nQ\otimes{}_nP\arrow[d,"{\alpha^n_{Q,P}}"]\\
\mathcal Tor_1(P,Q)\arrow[r,"\sim"]&\mathcal Tor_1(Q,P).
\end{tikzcd}\tag{2.1.10.1}\]
\end{SGAresult}
\SourcePage{240}
\begin{proof}
2.1.8的方法立即把我们归结到 $P=Q=(\mathbf Z/n\mathbf Z)_T$，并须证明：分别用（2.1.8.2）\Term{resolution}{解消} $\mathcal Tor_1$ 中第一个和第二个变元所得到的两个\Term{isomorphism}{同构}（2.1.8.1）相差一个负号（也就是说，$\mathcal Tor_1^{\mathbf Z}(\mathbf Z/n\mathbf Z,\mathbf Z/n\mathbf Z)$ 上的\Term{symmetry-isomorphism}{对称同构}是 $-\operatorname{id}$）。为此，显然可以假设 $T$ 是点\Term{topos}{拓扑斯}，即在通常的\Term{commutative-group}{交换群}\Term{category}{范畴}中工作。设
\[L_\bullet=(0\longrightarrow\mathbf Z\xrightarrow{n}\mathbf Z\longrightarrow0),\]
作为 $P=Q=\mathbf Z/n\mathbf Z$ 的\Term{flat-resolution}{平坦解消}。用 $L_\bullet\otimes L_\bullet$ 计算 $\mathcal Tor_1(P,Q)=\mathcal Tor_1(\mathbf Z/n\mathbf Z,\mathbf Z/n\mathbf Z)$，则它与 $H_1(L_\bullet\otimes\mathbf Z/n\mathbf Z)$（用来定义 $\alpha$）及 $H_1(\mathbf Z/n\mathbf Z\otimes L_\bullet)$（用来定义 $\beta$）的\Term{isomorphism}{同构}，分别由 $L_\bullet\otimes L_\bullet$ 的第二个和第一个因子的\Term{augmentation}{增广} $L_\bullet\longrightarrow\mathbf Z/n\mathbf Z$ 定义。立即可见，$\alpha(1,1)$ 由圈
\[-1_0\otimes1_1+1_1\otimes1_0\in(L_\bullet\otimes L_\bullet)_1
=L_0\otimes L_1+L_1\otimes L_0\]
表示（注意 $1_0,1_1$ 分别表示 $L_0,L_1$ 中的元素 $1$），而 $\beta(1,1)$ 由圈 $1_0\otimes1_1-1_1\otimes1_0$ 表示。这就证明了2.1.10。
\end{proof}
要结束关于\Term{pairing}{配对} $\alpha$ 的这些\Term{generic}{一般}讨论，还须说明 $n$ 变化时各\Term{homomorphism}{同态} $\alpha^n_{P,Q}$ 的关系。主要结果如下。
\begin{SGAresult}{\Term{lemma}{引理}}{2.1.11}
设 $P,Q$ 为 $T$ 的两个\Term{commutative-group}{交换群}，$n,n'>0$ 为整数，且 $n'$ 是 $n$ 的倍数，记 $n'=mn$。则下图交换：
\SourcePage{241}
\[\begin{tikzcd}[column sep=large]
&{}_{n'}P\otimes{}_nQ\arrow[dl,"\rho_m\otimes\operatorname{id}_{{}_nQ}"']\arrow[dr,"\operatorname{id}_{{}_{n'}P}\otimes i_m"]&\\
{}_nP\otimes{}_nQ\arrow[dr,"{\alpha^n_{P,Q}}"']&&{}_{n'}P\otimes{}_{n'}Q\arrow[dl,"{\alpha^{n'}_{P,Q}}"]\\
&\mathcal Tor_1(P,Q)&
\end{tikzcd}\tag{2.1.11.1}\]
其中 $\rho_m:{}_{n'}P\longrightarrow{}_nP$ 由 $m\operatorname{id}_P$ 诱导，$i_m:{}_nQ\longrightarrow{}_{n'}Q$ 为\Term{inclusion}{包含映射}。与（2.1.11.1）\Term{symmetric}{对称}的类似图（2.11.2）也交换：其顶点与（2.11.1）相同，只把上顶点 ${}_{n'}P\otimes{}_nQ$ 换为 ${}_nP\otimes{}_{n'}Q$。
\end{SGAresult}
\begin{proof}
只须证明第一个断言；由\Term{symmetric}{对称}性及2.1.10，便可得到第二个。用2.1.8的方法，可归结到 $T$ 是点\Term{topos}{拓扑斯}，即 $P,Q$ 是通常的\Term{commutative-group}{交换群}，且 $P=\mathbf Z/n'\mathbf Z,Q=\mathbf Z/n'\mathbf Z$。实际上，对 $P$ 只用到 $n'P=0$。借助 $Q$ 的\Term{resolution}{解消} $0\longrightarrow\mathbf Z\xrightarrow{n'}\mathbf Z\longrightarrow0$，将 $\mathcal Tor_1(P,Q)$ 等同于 ${}_{n'}P=P$；同样，借助 ${}_nQ\simeq\mathbf Z/n\mathbf Z$ 的\Term{resolution}{解消} $0\longrightarrow\mathbf Z\xrightarrow n\mathbf Z\longrightarrow0$，将 $\mathcal Tor_1({}_nP,{}_nQ)$ 等同于 ${}_nP$；这里\Term{augmentation}{增广}把 $1_0$ 映到 $m1_Q$（$1_Q$ 表示 $Q=\mathbf Z/n'\mathbf Z$ 中的元素 $1$）。由\Term{inclusion}{包含映射} ${}_nP\hookrightarrow P$、${}_nQ\hookrightarrow Q$ \Term{derive}{导出}的\Term{homomorphism}{同态}
\[\begin{tikzcd}
\mathcal Tor_1({}_nP,{}_nQ)\arrow[r]&\mathcal Tor_1(P,Q)\\
{}_nP\arrow[u,"\wr"']&{}_{n'}P=P\arrow[u,"\wr"']
\end{tikzcd}\]
\SourcePage{242}
就是自然\Term{inclusion}{包含映射} ${}_nP\hookrightarrow{}_{n'}P$。因此，用 ${}_nQ$ 的\Term{generator}{生成元} $m1_Q$ 所给的等同 ${}_nQ\simeq\mathbf Z/n\mathbf Z$，把（2.1.11.1）的上顶点等同于 $P/nP$，该图变成
\[\begin{tikzcd}
&P/nP\arrow[dl,"{x\mapsto mx}"']\arrow[dr,"{x\mapsto mx}"]&\\
{}_nP\arrow[dr,"{x\mapsto x}"']&&P\arrow[dl,"{x\mapsto x}"]\\
&{}_{n'}P&
\end{tikzcd}\]
它显然交换。
\end{proof}
\begin{SGAresult}{推论}{2.1.12}
沿用2.1.11的记号，下图交换：
\[\begin{tikzcd}
{}_{n'}P\otimes{}_{n'}Q\arrow[r,"m\otimes m"]\arrow[d,"{\alpha^{n'}_{P,Q}}"']&{}_nP\otimes{}_nQ\arrow[d,"{\alpha^n_{P,Q}}"]\\
{}_{n'}\mathcal Tor_1(P,Q)\arrow[r,"m"]&{}_n\mathcal Tor_1(P,Q).
\end{tikzcd}\tag{2.1.12.1}\]
描述水平\Term{arrow}{箭头}时，$m$ 表示由 $m\operatorname{id}_E$ 诱导的\Term{homomorphism}{同态} ${}_{n'}E\longrightarrow{}_nE$（$E$ 分别取 $P,Q,\mathcal Tor_1(P,Q)$）。
\end{SGAresult}
\begin{proof}
由2.1.11，对 $x\in{}_{n'}P(S)$、$y\in{}_{n'}Q(S)$ 有
$\alpha^n(mx\otimes my)=\alpha^{n'}(x\otimes my)$；又因 $x\otimes my=m(x\otimes y)$，右边等于 $m\alpha^{n'}(x\otimes y)$。这就得到（2.1.12.1）的\Term{commutativity}{交换性}。
\end{proof}
\SourcePage{243}
\textbf{2.1.13.} 设 $\ell$ 为\Term{prime-number}{素数}。对每个\Term{commutative-group}{交换群} $E$，用 $T_\ell(E)$ 表示\Term{projective-system}{投射系}
\[T_\ell(E)=({}_{\ell^\nu}E)_{\nu\geq0},\tag{2.1.13.1}\]
它显然\Term{functoriality}{函子性}地依赖于 $E$。对形如 $\ell^\nu,\ell^{\nu'}$（$\nu'\geq\nu$）的 $n,n'$，2.1.12的相容性表明，随 $\nu$ 变化的 $\alpha^{\ell^\nu}_{P,Q}$ 定义了一个对 $P,Q$ 都是\Term{functoriality}{函子性}的\Term{projective-system}{投射系}\Term{homomorphism}{同态}
\[\alpha^{(\ell)}_{P,Q}:T_\ell(P)\otimes T_\ell(Q)\longrightarrow T_\ell(\mathcal Tor_1(P,Q)).\tag{2.1.13.2}\]
知道每个\Term{prime-number}{素数} $\ell$ 对应的这些\Term{homomorphism}{同态}，显然\Term{equivalence}{等价}于知道每个整数 $n$ 对应的 $\alpha^n_{P,Q}$。\Term{antisymmetric}{反对称}关系2.1.10于是变成每个 $\ell$ 对应的类似 $\ell$-进\Term{antisymmetric}{反对称}关系，由下图的\Term{anticommutativity}{反交换性}表示：
\[\begin{tikzcd}
T_\ell(P)\otimes T_\ell(Q)\arrow[r]\arrow[d,"{\alpha^{(\ell)}_{P,Q}}"']&T_\ell(Q)\otimes T_\ell(P)\arrow[d,"{\alpha^{(\ell)}_{Q,P}}"]\\
T_\ell(\mathcal Tor_1(P,Q))\arrow[r]&T_\ell(\mathcal Tor_1(Q,P)).
\end{tikzcd}\tag{2.1.13.3}\]
水平\Term{arrow}{箭头}是\Term{symmetric}{对称}\Term{arrow}{箭头}。

\textbf{注记 2.1.14.} 可以把\Term{homomorphism}{同态}（2.1.1）的构造及关于它的结果2.1.10、2.1.11推广到以下情形：给定 $T$ 上的一个交换\Term{ring}{环} $A$ 及 $A$ 的两个\Term{ideal}{理想} $\mathfrak a,\mathfrak b$（在前面处理的情形，它们都是 $T$ 上\Term{constant}{常值}\Term{ring}{环} $\mathbf Z_T$ 的同一个\Term{ideal}{理想} $n\mathbf Z_T$）。
\SourcePage{244}
此时，定义一个对 $A$-\Term{module}{模} $P,Q$ 都是\Term{functoriality}{函子性}的\Term{homomorphism}{同态}
\[{}_{\mathfrak a}P\otimes{}_{\mathfrak b}Q\xrightarrow{\alpha}
\mathcal Hom_A(\mathfrak a\cap\mathfrak b/\mathfrak a\mathfrak b,\mathcal Tor_1^A(P,Q))\tag{2.1.14.1}\]
为如下\Term{homomorphism}{同态}的复合：
\[\begin{split}
{}_{\mathfrak a}P\otimes{}_{\mathfrak b}Q&\simeq
\mathcal Hom(A/\mathfrak a,P)\otimes\mathcal Hom(A/\mathfrak b,Q)\\
&\longrightarrow\mathcal Hom(\mathcal Tor_1(A/\mathfrak a,A/\mathfrak b),\mathcal Tor_1(P,Q)),
\end{split}\]
这里用到熟知的\Term{canonical-isomorphism}{典范同构}
\[\mathcal Tor_1(A/\mathfrak a,A/\mathfrak b)\xrightarrow{\sim}
\mathfrak a\cap\mathfrak b/\mathfrak a\mathfrak b.\tag{2.1.14.2}\]
事实上，2.1.10中的\Term{signature}{符号}就“隐藏”在这个\Term{isomorphism}{同构}的选择中，取决于我们使用 $A/\mathfrak a$ 的\Term{resolution}{解消} $0\longrightarrow\mathfrak a\longrightarrow A\longrightarrow0$，还是 $A/\mathfrak b$ 的\Term{resolution}{解消} $0\longrightarrow\mathfrak b\longrightarrow A\longrightarrow0$。2.1.10在这里推广为下图的\Term{anticommutativity}{反交换性}：
\[\resizebox{\linewidth}{!}{\begin{tikzcd}[ampersand replacement=\&,column sep=small]
{}_{\mathfrak a}P\otimes{}_{\mathfrak b}Q\arrow[r]\arrow[d]\&{}_{\mathfrak b}Q\otimes{}_{\mathfrak a}P\arrow[d]\\
\mathcal Hom(\mathfrak a\cap\mathfrak b/\mathfrak a\mathfrak b,\mathcal Tor_1(P,Q))\arrow[r]\&
\mathcal Hom(\mathfrak a\cap\mathfrak b/\mathfrak a\mathfrak b,\mathcal Tor_1(Q,P)).
\end{tikzcd}}\tag{2.1.14.3}\]
水平\Term{arrow}{箭头}是\Term{symmetric}{对称}\Term{arrow}{箭头}。当 $\mathfrak a=\mathfrak b$ 时，这种\Term{anticommutativity}{反交换性}也可表述为：$\mathcal Tor_1(A/\mathfrak a,A/\mathfrak a)$ 中的\Term{symmetric}{对称}\Term{automorphism}{自同构}等于 $-\operatorname{id}$。特别地，若 $\mathfrak a=\mathfrak b=nA$，其中 $n$ 是 $A$ 的\Term{section}{截面}，则乘以 $n$ 诱导\Term{epimorphism}{满态射}
\[A/\mathfrak a\longrightarrow\mathfrak a\cap\mathfrak a/\mathfrak a\mathfrak a=\mathfrak a/\mathfrak a^2,\]
从而有一个对 $A$-\Term{module}{模} $T$ 是\Term{functoriality}{函子性}的\Term{injective-homomorphism}{单同态}
\SourcePage{245}
\[\mathcal Hom(\mathfrak a/\mathfrak a^2,T)\hookrightarrow{}_{\mathfrak a}T.\]
因而此时（2.1.14.1）\Term{solvable}{可解}释为更接近（2.1.1）\Term{form}{形式}的\Term{homomorphism}{同态}
\[{}_nP\otimes{}_nQ\longrightarrow{}_n\mathcal Tor_1(P,Q).\tag{2.1.14.4}\]
仍假设 $\mathfrak a=\mathfrak b=nA$，并假设 $n$ 是 $A$ 的\Term{regular-section}{正则截面}，则\Term{spectral-sequence}{谱序列}（2.1.4）在当前情形仍有意义，其低次项\Term{exact-sequence}{正合列}推广了（2.1.5）（现在 $Q$ 是被 $\mathfrak a$ \Term{annihilate}{零化}的 $A$-\Term{module}{模}）：
\[\begin{split}
\mathcal Tor_2^A(P,Q)&\longrightarrow\mathcal Tor_2^{A/nA}(P_n,Q)
\longrightarrow{}_nP\otimes Q\xrightarrow{\alpha}\mathcal Tor_1^A(P,Q)\\
&\longrightarrow\mathcal Tor_1^{A/nA}(P_n,Q)\longrightarrow0.
\end{split}\tag{2.1.14.5}\]
回到 $\mathfrak a,\mathfrak b$ 任意的情形，并给定 $A$ 的两个\Term{section}{截面} $p,q$，则有如下\Term{commutative-diagram}{交换图}（（2.1.11.1）本质上是其中 $q=1$ 的特例）：
\[\resizebox{\linewidth}{!}{\begin{tikzcd}[ampersand replacement=\&,column sep=small,row sep=large]
\&{}_{p\mathfrak a}P\otimes{}_{q\mathfrak b}Q\arrow[dl,"p\otimes q"']\arrow[dr,"i\otimes j"]\&\\
{}_{\mathfrak a}P\otimes{}_{\mathfrak b}Q\arrow[d,"\alpha"']\&\&{}_{pq\mathfrak a}P\otimes{}_{pq\mathfrak b}Q\arrow[d,"\alpha"]\\
\mathcal Hom(\mathfrak a\cap\mathfrak b/\mathfrak a\mathfrak b,\mathcal Tor_1(P,Q))\arrow[dr,"\mathrm{incl.}"']\&\&
\mathcal Hom(pq\mathfrak a\cap pq\mathfrak b/(pq\mathfrak a)(pq\mathfrak b),\mathcal Tor_1(P,Q))\arrow[dl,"(pq)^*"]\\
\&\mathcal Hom(\mathfrak a\cap\mathfrak b/pq\mathfrak a\mathfrak b,\mathcal Tor_1(P,Q))\&
\end{tikzcd}}\tag{2.1.14.6}\]
\SourcePage{246}
其中 $i:{}_{p\mathfrak a}P\longrightarrow{}_{pq\mathfrak a}P$、$j:{}_{q\mathfrak b}Q\longrightarrow{}_{pq\mathfrak b}Q$ 是\Term{canonical}{典范}\Term{inclusion}{包含映射}；$(pq)^*$ 是由乘以 $pq$ 所诱导的\Term{surjective-homomorphism}{满同态}
\[\mathfrak a\cap\mathfrak b/pq\mathfrak a\mathfrak b\longrightarrow
pq\mathfrak a\cap pq\mathfrak b/(pq\mathfrak a)(pq\mathfrak b)\]
\Term{action}{作用} $\mathcal Hom$ 而得到的\Term{injective-homomorphism}{单同态}。当 $\mathfrak a=\mathfrak b=nA$ 时，该图可写成更简单的\Term{form}{形式}（$q=1$ 的情形已在（2.1.11.1）中给出）。

\textbf{2.2.} 现在除了 $P,Q$，再考虑一个\Term{commutative-group}{交换群} $G$。我们已在（1.1.3）中定义了一个对 $P,Q,G$ 都是\Term{functoriality}{函子性}的\Term{canonical}{典范}\Term{homomorphism}{同态}
\[\operatorname{Biext}^1(P,Q;G)\longrightarrow\operatorname{Hom}(\mathcal Tor_1(P,Q),G).\tag{2.2.1}\]
另一方面，若 $n>0$ 为整数（相应地，若 $\ell$ 为\Term{prime-number}{素数}），则（2.1.1）的\Term{homomorphism}{同态} $\alpha^n_{P,Q}$（相应地，（2.1.13.2）的\Term{homomorphism}{同态} $\alpha^{(\ell)}_{P,Q}$）给出\Term{homomorphism}{同态}
\[\operatorname{Hom}(\mathcal Tor_1(P,Q),G)\longrightarrow\operatorname{Hom}({}_nP\otimes{}_nQ,G)\]
（相应地，给出\Term{homomorphism}{同态}
\[\operatorname{Hom}(\mathcal Tor_1(P,Q),G)\longrightarrow
\operatorname{Hom}(T_\ell(P)\otimes T_\ell(Q),T_\ell(G))\]
）。这些\Term{homomorphism}{同态}同样对 $P,Q,G$ 都是\Term{functoriality}{函子性}的。将它们与（2.2.1）复合，得到对 $P,Q,G$ 都是\Term{functoriality}{函子性}的\Term{canonical}{典范}\Term{homomorphism}{同态}
\[\varphi^n=\varphi^n_{P,Q}:\operatorname{Biext}^1(P,Q;G)\longrightarrow
\operatorname{Hom}({}_nP\otimes{}_nQ,G),\tag{2.2.2}\]
\[\varphi^{(\ell)}=\varphi^{(\ell)}_{P,Q}:\operatorname{Biext}^1(P,Q;G)\longrightarrow
\operatorname{Hom}(T_\ell(P)\otimes T_\ell(Q),T_\ell(G)).\tag{2.2.3}\]
对
\[\xi\in\operatorname{Biext}^1(P,Q;G),\tag{2.2.4}\]
\SourcePage{247}
用 $\varphi_\xi^n$（相应地 $\varphi_\xi^{(\ell)}$）表示它在（2.2.2）（相应地（2.2.3））下的像：
\[\varphi_\xi^n:{}_nP\otimes{}_nQ\longrightarrow G\quad\text{（也可取值于 }{}_nG\text{）},\tag{2.2.5}\]
\[\varphi_\xi^{(\ell)}:T_\ell(P)\otimes T_\ell(Q)\longrightarrow T_\ell(G).\tag{2.2.6}\]

\textbf{2.2.7.} 显然，知道 $\varphi_\xi^{(\ell)}$ \Term{equivalence}{等价}于知道所有 $n=\ell^\nu$（$\nu\geq0$）对应的 $\varphi_\xi^{(n)}$；固定 $n$ 时，知道 $\varphi_\xi^{(n)}$ \Term{equivalence}{等价}于知道\Term{homomorphism}{同态}
\[\mathcal Tor_1(P,Q)\longrightarrow G\tag{2.2.8}\]
在 $\alpha^n_{P,Q}:{}_nP\otimes{}_nQ\longrightarrow\mathcal Tor_1(P,Q;G)$ 的像上的\Term{restriction}{限制}，这个\Term{homomorphism}{同态}就是 $\xi$ 在（2.2.1）下的像。注意，当 $P=nP,Q=nQ$ 时，上述像等于 ${}_n\mathcal Tor_1(P,Q)$；这容易由\Term{exact-sequence}{正合列}（1.3.5）及 $Q$ 对应的类似\Term{exact-sequence}{正合列}得到。另一方面，若 $P,Q$ 都\Term{divisible}{可除}，即对每个整数 $n>0$ 都有 $nP=P,nQ=Q$，则容易\Term{kernel}{核}实，$\mathcal Tor_1(P,Q)$ 是\Term{torsion}{挠的}，即
\[\mathcal Tor_1(P,Q)=\varinjlim_n{}_n\mathcal Tor_1(P,Q).\]
所以，此时知道所有 $n$ 对应的 $\varphi_\xi^n$ \Term{equivalence}{等价}于知道（2.2.8）；而对给定\Term{prime-number}{素数} $\ell$，知道 $\varphi_\xi^{(\ell)}$ \Term{equivalence}{等价}于知道（2.2.8）在 $\mathcal Tor_1(P,Q)$ 的 $\ell$-\Term{primary-component}{准素分量}上的\Term{restriction}{限制}。现在具体说明与\Term{biextension}{双扩张}关联的\Term{pairing}{配对}（2.2.4）、（2.2.6）的\Term{symmetric}{对称}性质。先注意有\Term{commutative-diagram}{交换图}
\[\begin{tikzcd}
\operatorname{Biext}^1(P,Q;G)\arrow[r,"\sim"]\arrow[d]&\operatorname{Biext}^1(Q,P;G)\arrow[d]\\
\operatorname{Hom}(\mathcal Tor_1(P,Q),G)\arrow[r,"\sim"]&\operatorname{Hom}(\mathcal Tor_1(Q,P),G).
\end{tikzcd}\tag{2.2.9}\]
\SourcePage{248}
水平\Term{arrow}{箭头}是\Term{symmetric}{对称}\Term{arrow}{箭头}（第一个在第VII讲2.7中定义），竖直\Term{arrow}{箭头}是（2.2.1）。\Term{commutativity}{交换性}事实上来自下面更完整的图中两个方块各自的\Term{commutativity}{交换性}：
\[\begin{tikzcd}
\operatorname{Biext}^1(P,Q;G)\arrow[r,"\sim"]\arrow[d]&\operatorname{Biext}^1(Q,P;G)\arrow[d]\\
\operatorname{Ext}^1(P\mathbin{\overset{\mathbf L}{\otimes}}Q,G)\arrow[r,"\sim"]\arrow[d]&
\operatorname{Ext}^1(Q\mathbin{\overset{\mathbf L}{\otimes}}P,G)\arrow[d]\\
\operatorname{Hom}(\mathcal Tor_1(P,Q),G)\arrow[r,"\sim"]&\operatorname{Hom}(\mathcal Tor_1(Q,P),G).
\end{tikzcd}\tag{2.2.10}\]
每列中的第一个竖直\Term{arrow}{箭头}是第VII讲3.6.5的\Term{canonical-isomorphism}{典范同构}，第二个是显然的\Term{homomorphism}{同态}。（2.2.10）中第二个方块的\Term{commutativity}{交换性}是平凡的：按定义，\Term{symmetry-isomorphism}{对称同构}
$P\mathbin{\overset{\mathbf L}{\otimes}}Q\xrightarrow{\sim}Q\mathbin{\overset{\mathbf L}{\otimes}}P$
诱导\Term{symmetry-isomorphism}{对称同构} $\mathcal Tor_1(P,Q)\xrightarrow{\sim}\mathcal Tor_1(Q,P)$；而对变动的\Term{complex}{复形} $L_\bullet$，\Term{canonical}{典范}\Term{homomorphism}{同态} $\operatorname{Ext}^1(L_\bullet,G)\longrightarrow\operatorname{Hom}(H_1(L_\bullet),G)$ 对 $L_\bullet$ 是\Term{functoriality}{函子性}的。至于第一个方块的\Term{commutativity}{交换性}，则来自同时写出\Term{symmetry-isomorphism}{对称同构}和第VII讲3.6.5定义的竖直\Term{isomorphism}{同构}的具体\Term{form}{形式}；\Term{kernel}{核}实的细节留给读者。这里要用到：\Term{symmetry-isomorphism}{对称同构}
$L_\bullet(P)\otimes L_\bullet(Q)\xrightarrow{\sim}L_\bullet(Q)\otimes L_\bullet(P)$ 在\Term{degree}{次数}1诱导\Term{isomorphism}{同构}
\[\begin{split}
L_0(P)\otimes L_1(Q)+L_1(P)\otimes L_0(Q)
\simeq L_0(Q)\otimes L_1(P)+L_1(Q)\otimes L_0(P),
\end{split}\]
它在相应\Term{direct-summand}{直和项}上由底层未\Term{graded-module}{分次模}的通常\Term{symmetry-isomorphism}{对称同构}给出
\SourcePage{249}
（因此没有\Term{signature}{符号}介入）。
考虑到（2.2.9）的\Term{commutativity}{交换性}，（2.1.10.1）的\Term{anticommutativity}{反交换性}用 $\varphi^n$ 或 $\varphi^{(\ell)}$ 表示，便得到以下关系。
\begin{SGAresult}{命题}{2.2.11}
以下两图\Term{anticommutativity}{反交换}，水平\Term{arrow}{箭头}是\Term{symmetric}{对称}\Term{arrow}{箭头}：
\[\begin{tikzcd}
\operatorname{Biext}^1(P,Q;G)\arrow[r,"\sim"]\arrow[d,"{\varphi^n_{P,Q}}"']&\operatorname{Biext}^1(Q,P;G)\arrow[d,"{\varphi^n_{Q,P}}"]\\
\operatorname{Hom}({}_nP\otimes{}_nQ,G)\arrow[r,"\sim"]&\operatorname{Hom}({}_nQ\otimes{}_nP,G),
\end{tikzcd}\tag{2.2.11.1}\]
\[\resizebox{\linewidth}{!}{\begin{tikzcd}[ampersand replacement=\&,column sep=small]
\operatorname{Biext}^1(P,Q;G)\arrow[r,"\sim"]\arrow[d,"{\varphi^{(\ell)}_{P,Q}}"']\&\operatorname{Biext}^1(Q,P;G)\arrow[d,"{\varphi^{(\ell)}_{Q,P}}"]\\
\operatorname{Hom}(T_\ell(P)\otimes T_\ell(Q),T_\ell(G))\arrow[r,"\sim"]\&
\operatorname{Hom}(T_\ell(Q)\otimes T_\ell(P),T_\ell(G)).
\end{tikzcd}}\tag{2.2.11.2}\]
\end{SGAresult}
指出以下重要推论。
\begin{SGAresult}{推论}{2.2.12}
假设 $P=Q$，且 $\xi$ 是 $(P,Q)$ 的一个以 $G$ 为\Term{kernel}{核}的\Term{symmetric}{对称}（相应地，\Term{antisymmetric}{反对称}）\Term{biextension}{双扩张}的类，即 $s\xi=\xi$（相应地 $s\xi=-\xi$），其中 $s\xi$ 表示 $\xi$ 经\Term{symmetric}{对称}所得的像。则\Term{bilinear}{双线性}映射 $\varphi_\xi^n$（定义在 ${}_nP\times{}_nP$ 上）与 $\varphi_\xi^{(\ell)}$（定义在 $T_\ell(P)\times T_\ell(P)$ 上）是\Term{antisymmetric}{反对称}的（相应地，\Term{symmetric}{对称}的）。
\end{SGAresult}
\SourcePage{250}
\textbf{注记 2.2.13.} 假设 $\xi$ \Term{symmetric}{对称}。若 $n$（相应地 $\ell$）为奇数，则由2.2.12知，$\varphi_\xi^n$（相应地 $\varphi_\xi^{(\ell)}$）甚至是\Term{alternating-form}{交替形式}。若不\Term{restriction}{限制} $n$（相应地 $\ell$）为奇数，这个结论\Term{generic}{一般}不成立；$P$ 是通常的\Term{group}{群} $\mathbf Z/2\mathbf Z$ 且 $G=\mathbf Z/2\mathbf Z$ 的情形就已说明这一点。不过，假设 $2P=P$（特别地，$P$ 是\Term{abelian-scheme}{阿贝尔概形}且照常使用 fppf 或 fpqc \Term{topology}{拓扑}时如此），则上述\Term{form}{形式}仍为\Term{alternating-form}{交替的}。只须证明 $\varphi^{2^\nu}(x,x)=0$。由于 $2P=P$，可以假设给定的 $x\in{}_{2^\nu}P(S)$ 形如 $2y$，其中 $y\in{}_{2^{\nu+1}}P(S)$，于是
\[\varphi^{2^\nu}(x,x)=2\varphi^{2^{\nu+1}}(y,y)=0,\]
因为 $\varphi^{2^{\nu+1}}$ \Term{antisymmetric}{反对称}。

\textbf{2.3. 与“Cartier \Term{duality}{对偶性}”的联系。} 2.2借助\Term{isomorphism}{同构}\Term{theorem}{定理}第VII讲3.6.5及\Term{homomorphism}{同态}（2.2.1），描述了一个\Term{biextension}{双扩张}所定义的\Term{pairing}{配对}；这显然不是过去在\Term{abelian-scheme}{阿贝尔概形}情形采用的定义（那时甚至还没有\Term{biextension}{双扩张}这个概念！）。原来的定义通过 Cartier \Term{duality}{对偶性}进行；现在说明与这种\Term{duality}{对偶性}的关系。
重新考虑\Term{homomorphism}{同态}
\[\operatorname{Biext}^1(P,Q;G)\longrightarrow\operatorname{Hom}(P,\mathcal Ext^1(Q,G)),\tag{2.3.1}\]
它来自\Term{exact-sequence}{正合列}（1.4），而且在
\[\mathcal Hom(Q,G)=0\tag{2.3.2}\]
\SourcePage{251}
时是\Term{isomorphism}{同构}；特别地，当 $Q$ 是\Term{scheme}{概形} $S$ 上的\Term{abelian-scheme}{阿贝尔概形}，$G$ 是\Term{multiplicative-group}{乘法群} $\mathbf G_m$ 时，这个条件成立（第VII讲1.3.8）。另一方面，有显然的“\Term{restriction}{限制}到 $nQ$”\Term{homomorphism}{同态} $\mathcal Ext^1(Q,G)\longrightarrow\mathcal Ext^1(nQ,G)$，从而有复合\Term{homomorphism}{同态}
\[\begin{split}
\operatorname{Hom}(P,\mathcal Ext^1(Q,G))&\longrightarrow\operatorname{Hom}(P,\mathcal Ext^1(nQ,G))\\
&\longrightarrow\operatorname{Hom}({}_nP,{}_n\mathcal Ext^1(nQ,G)),
\end{split}\tag{2.3.4}\]
第二个\Term{arrow}{箭头}是\Term{restriction}{限制}到 ${}_nP$ 的\Term{homomorphism}{同态}。与（2.3.1）复合得到
\[\operatorname{Biext}^1(P,Q;G)\longrightarrow\operatorname{Hom}({}_nP,{}_n\mathcal Ext^1(nQ,G)).\tag{2.3.5}\]
另一方面，\Term{exact-sequence}{正合列}
\[0\longrightarrow{}_nQ\longrightarrow Q\xrightarrow{n}nQ\longrightarrow0\tag{2.3.6}\]
给出 $\mathcal Ext^i(-,G)$ 的\Term{exact-sequence}{正合列}，进而得到\Term{short-exact-sequence}{短正合列}
\[0\longrightarrow\mathcal Hom(Q,G)_n\longrightarrow\mathcal Hom({}_nQ,G)
\xrightarrow{\partial}{}_n\mathcal Ext^1(nQ,G)\longrightarrow0,\tag{2.3.7}\]
其中照常以 $E_n$ 表示 $E/nE$。我们要用的是\Term{canonical}{典范}\Term{surjective-homomorphism}{满同态}
\[\mathcal Hom({}_nQ,G)\xrightarrow{\partial}{}_n\mathcal Ext^1(nQ,G),\tag{2.3.8}\]
当 $\mathcal Hom(Q,G)_n=0$ 时它还是\Term{isomorphism}{同构}，特别地，（2.3.2）成立时如此。由（2.3.8）\Term{derive}{导出}\Term{homomorphism}{同态}
\[\begin{split}
\operatorname{Hom}({}_nP,\mathcal Hom({}_nQ,G))&\simeq\operatorname{Hom}({}_nP\otimes{}_nQ,G)\\
&\longrightarrow\operatorname{Hom}({}_nP,{}_n\mathcal Ext^1(nQ,G)),
\end{split}\tag{2.3.9}\]
当（2.3.2）成立时它是\Term{isomorphism}{同构}。所有上述\Term{homomorphism}{同态}对 $P,Q,G$ 都是\Term{functoriality}{函子性}的。\Term{homomorphism}{同态}（2.3.5）、（2.3.9）、（2.2.2）给出图
\SourcePage{252}
\[\begin{tikzcd}[column sep=large]
&\operatorname{Biext}^1(P,Q;G)\arrow[dl,"\varphi^n"']\arrow[dr,"\text{（2.3.5）}"]&\\
\operatorname{Hom}({}_nP\otimes{}_nQ,G)\arrow[rr,"\text{（2.3.9）}"']&&
\operatorname{Hom}({}_nP,{}_n\mathcal Ext^1(nQ,G)),
\end{tikzcd}\tag{2.3.10}\]
若（2.3.2）成立，水平\Term{arrow}{箭头}（2.3.9）就是\Term{isomorphism}{同构}。
\begin{SGAresult}{命题}{2.3.11}
图（2.3.10）\Term{anticommutativity}{反交换}。
\end{SGAresult}
当 $\mathcal Hom(Q,G)=0$ 时，该命题用\Term{homomorphism}{同态}（2.3.5）、（2.3.9）重新刻画了（2.2.2）的\Term{homomorphism}{同态} $\varphi^n$；这正是过去\Term{abelian-scheme}{阿贝尔概形}理论中采用的刻画（而且在那里有 $nQ=Q$）。
\begin{proof}[2.5.11的证明]
\footnote{更优美的证明见下文2.3.19。}
设 $\xi\in\operatorname{Biext}^1(P,Q;G)$，$\varphi_\xi^n:{}_nP\otimes{}_nQ\longrightarrow G$ 为关联\Term{pairing}{配对}，$u_\xi:{}_nP\longrightarrow\mathcal Ext^1(nQ,G)$ 为 $\xi$ 在（2.3.5）下的像。须证明 $-u_\xi$ 由 $\varphi_\xi$ 经（2.3.9）\Term{derive}{导出}，即为复合
\[{}_nP\xrightarrow{\varphi_\xi}\mathcal Hom({}_nQ,G)
\xrightarrow{\text{（2.3.8）}}{}_n\mathcal Ext^1(nQ,G).\tag{*}\]
为此，可以给定 ${}_nP$ 的一个\Term{section}{截面} $x$，证明它在 $-u_\xi$ 下的像等于在复合（*）下的像。（如果愿意，这样可归结到 $P=(\mathbf Z/n\mathbf Z)_T$，但我们不需要这一步。）

 a）分别选择 $P,Q$ 的\Term{flat-resolution}{平坦解消} $L_\bullet,M_\bullet$：
\[0\longrightarrow L_1\longrightarrow L_0\longrightarrow0,\qquad
0\longrightarrow M_1\longrightarrow M_0\longrightarrow0.\]
例如，$P$ 的\Term{flat-resolution}{平坦解消}可如下取得：
\SourcePage{253}
将 $P$ 写为\Term{flat}{平坦} $\mathbf Z_T$-\Term{module}{模} $L_0$ 的\Term{quotient}{商}，令 $L_1$ 为 $L_0\longrightarrow P$ 的\Term{kernel}{核}（它\Term{torsion-free}{无挠}，故\Term{flat}{平坦}）。由此得到\Term{complex}{复形} $L_\bullet\otimes M_\bullet$
\[0\longrightarrow L_1\otimes M_1\xrightarrow{d_2}
L_1\otimes M_0+L_0\otimes M_1\xrightarrow{d_1}L_0\otimes M_0\longrightarrow0,\]
我们用它的\Term{homology-sheaf}{同调层} $H_1$ 作为 $\mathcal Tor_1(P,Q)$。以 $L'_0$（相应地 $M'_0$）表示 $L_0$（相应地 $M_0$）中 ${}_nP$（相应地 ${}_nQ$）的\Term{inverse-image}{逆像}\Term{subsheaf}{子层}；于是\Term{complex}{复形} $0\longrightarrow L_1\longrightarrow L'_0\longrightarrow0$（相应地 $0\longrightarrow M_1\longrightarrow M'_0\longrightarrow0$）是 ${}_nP$（相应地 ${}_nQ$）的\Term{flat-resolution}{平坦解消}。此时\Term{homomorphism}{同态}
\[{}_nP\otimes{}_nQ\longrightarrow\mathcal Tor_1(P,Q)=H_1(L_\bullet\otimes M_\bullet)\]
可以由\Term{homomorphism}{同态}
\[L'_0\otimes M'_0\longrightarrow H_1(L_\bullet\otimes M_\bullet),\qquad
x_0\otimes y_0\longmapsto\text{1-圈 }[nx_0]\otimes y_0-x_0\otimes[ny_0]\text{ 的类}\]
通过取\Term{quotient}{商}具体给出。方\Term{lie-bracket}{括号}表示把有关元素看成 $L_0$ 的\Term{subsheaf}{子层} $L_1$（相应地 $M_0$ 的\Term{subsheaf}{子层} $M_1$）的\Term{section}{截面}。这个具体描述来自2.1.10证明中的计算。

 b）取 $G$ 的\Term{injective-resolution}{内射解消}
\[0\longrightarrow C^0\xrightarrow{d^0}C^1\longrightarrow\cdots,\]
并作等同
\[\operatorname{Biext}^1(P,Q;G)\simeq\operatorname{Ext}^1(P\mathbin{\overset{\mathbf L}{\otimes}}Q,G)
=H^1(\operatorname{Hom}^\bullet(L_\bullet\otimes M_\bullet,C^\bullet)).\]
于是元素 $\xi$ 可以解释为一对\Term{homomorphism}{同态} $f_1,f_0$，组成\Term{anticommutativity}{反交换}图
\SourcePage{254}
\[\begin{tikzcd}
L_1\otimes M_0+L_0\otimes M_1\arrow[r,"d_1"]\arrow[d,"f_1"']&L_0\otimes M_0\arrow[d,"f_0"]\\
C^0\arrow[r,"d^0"]&C^1.
\end{tikzcd}\tag{2.3.12}\]
按定义，\Term{pairing}{配对} $\varphi_\xi^n:{}_nP\otimes{}_nQ\longrightarrow G$ 由\Term{homomorphism}{同态} $\varphi:L'_0\otimes M'_0\longrightarrow G$ 通过取\Term{quotient}{商}得到，其中
\[\varphi(x_0\otimes y_0)=f_1([nx_0]\otimes y_0-x_0\otimes[ny_0])\in G=\ker(C^0\longrightarrow C^1).\tag{2.3.13}\]
可以假设给定的 ${}_nP$ 的\Term{section}{截面} $x$ \Term{lifting}{提升}为 $L'_0$ 的\Term{section}{截面} $x_0$（因为断言（2.3.11）在 $T$ 上是\Term{local-scheme}{局部}的）。把上式中的 $y_0$ 看成变量，就定义\Term{homomorphism}{同态}
\[\varphi':M'_0\longrightarrow G,\qquad
y_0\longmapsto f_1([nx_0]\otimes y_0-x_0\otimes[ny_0]),\tag{2.3.14}\]
取\Term{quotient}{商}便得到与 $\varphi_\xi$ 和 $x$ 关联的\Term{homomorphism}{同态}
\[{}_nQ\longrightarrow G,\qquad y\longmapsto\varphi_\xi(x,y).\]

 c）须求这个\Term{homomorphism}{同态}在与\Term{exact-sequence}{正合列}（2.3.6）关联的（2.3.7）的\Term{operator}{算子} $\partial$ 下的像。显然，$nQ\simeq Q/{}_nQ$ 也\Term{isomorphism}{同构}于 $M_0/M'_0$，所以有\Term{exact-sequence}{正合列}之间的\Term{homomorphism}{同态}
\SourcePage{255}
\[\begin{tikzcd}
0\arrow[r]&M'_0\arrow[r]\arrow[d]&M_0\arrow[r]\arrow[d]&M_0/M'_0\arrow[r]\arrow[d,"\wr"]&0\\
0\arrow[r]&{}_nQ\arrow[r]&Q\arrow[r]&nQ\arrow[r]&0,
\end{tikzcd}\]
最后一个竖直\Term{arrow}{箭头}为\Term{isomorphism}{同构}。因此，所求的 $\mathcal Ext^1(nQ,G)$ 的\Term{section}{截面}也可计算为 $\varphi'$（2.3.14）在\Term{coboundary-operator}{余边算子}
\[\mathcal Hom(M'_0,G)\xrightarrow{\partial}\mathcal Ext^1(M_0/M'_0,G)\]
下的像；或者，它是第二个成员中由如下全局 $\operatorname{Ext}^1$ 的元素所定义的\Term{section}{截面}：$\varphi'$ 在\Term{coboundary-operator}{余边算子}
\[\operatorname{Hom}(M'_0,G)\longrightarrow\operatorname{Ext}^1(M_0/M'_0,G)\]
下的像。按照\Term{coboundary-operator}{余边算子}的通常构造，可以这样计算：将复合（2.3.13）
\[M'_0\xrightarrow{\varphi'}G\longrightarrow C^0\]
延拓为\Term{homomorphism}{同态}
\[\psi:M_0\longrightarrow C^0;\]
这是可能的，因为 $C^0$ \Term{injective-object}{内射}。再考虑使下图交换的唯一\Term{homomorphism}{同态}
\[\lambda:M_0/M'_0\longrightarrow C^1:\tag{2.3.15}\]
\[\begin{tikzcd}
0\arrow[r]&M'_0\arrow[r]\arrow[d,"\varphi'"']&M_0\arrow[r]\arrow[d,"\psi"']&M_0/M'_0\arrow[r]\arrow[d,"\lambda"]&0\\
0\arrow[r]&G\arrow[r]&C^0\arrow[r]&C^1.&
\end{tikzcd}\]
\SourcePage{256}
显然，（2.3.15）的值都是1-\Term{cocycle}{余圈}，即它是一个1-\Term{cocycle}{余圈}；它代表所求的 $\mathcal Ext^1(nQ,G)=\mathcal Ext^1(M_0/M'_0,G)$ 的元素。

 d）还须具体写出由 $\xi$ 和 $x$ 经图（2.3.10）中的\Term{arrow}{箭头}（2.3.5）\Term{derive}{导出}的 $\mathcal Ext^1(nQ,G)$ 的\Term{section}{截面}。为此，可以暂不使用 $P$ 的\Term{section}{截面} $x$ 被 $n$ \Term{annihilate}{零化}这一事实；由选定的 $x$ 在 $L_0$ 中的\Term{lifting}{提升} $x_0$，可作如下描述。它定义\Term{augmentation-morphism}{增广同态}
\[u:M_\bullet\longrightarrow L_\bullet\otimes M_\bullet,
\qquad y_0\longmapsto x_0\otimes y_0,\quad y_1\longmapsto x_0\otimes y_1.\]
由（2.3.1）\Term{derive}{导出}的 $\mathcal Ext^1(Q,G)$ 的\Term{section}{截面}，是由 $\operatorname{Ext}^1(L_\bullet\otimes M_\bullet,G)$ 的元素 $\xi$ 与上述\Term{homomorphism}{同态}复合得到的\Term{section}{截面}；因此，它由如下\Term{degree}{次数}1的\Term{homomorphism}{同态}给出：
\[\begin{tikzcd}
0\arrow[r]&M_1\arrow[r,"d_1"]\arrow[d,"{y_1\mapsto f_1(x_0\otimes y_1)}"']&M_0\arrow[r]\arrow[d,"{y_0\mapsto f_1(x_0\otimes y_0)}"]&0\\
0\arrow[r]&C^0\arrow[r,"d^0"]&C^1\arrow[r]&\cdots .
\end{tikzcd}\]
这个方块\Term{anticommutativity}{反交换}。为把此类\Term{restriction}{限制}到 $nQ\hookrightarrow Q$，使用\Term{isomorphism}{同构} $nQ\simeq M_0/M'_0$；它给出 $nQ$ 的\Term{flat-resolution}{平坦解消} $0\longrightarrow M'_0\longrightarrow M_0\longrightarrow0$，而\Term{inclusion}{包含映射} $nQ\hookrightarrow Q$ \Term{lifting}{提升}为\Term{resolution}{解消}之间的\Term{homomorphism}{同态}
\[\begin{tikzcd}
0\arrow[r]&M'_0\arrow[r]\arrow[d,"{y'_0\mapsto[ny'_0]}"']&M_0\arrow[r]\arrow[d,"{y_0\mapsto ny_0}"]&0\\
0\arrow[r]&M_1\arrow[r]&M_0\arrow[r]&0.
\end{tikzcd}\]
\SourcePage{257}
复合后便看出，所求的 $\mathcal Ext^1(nQ,G)$ 的\Term{section}{截面}关联于 $\operatorname{Ext}^1(nQ,G)$ 中由如下\Term{degree}{次数}1的\Term{complex-morphism}{复形同态}代表的元素：
\[\begin{tikzcd}
0\arrow[r]&M'_0\arrow[r,"d_1"]\arrow[d,"g_1"']&M_0\arrow[r]\arrow[d,"g_0"]&0\\
0\arrow[r]&C^0\arrow[r]&C^1\arrow[r]&\cdots,
\end{tikzcd}\]
其中
\[g_1(y'_0)=f_1(x_0\otimes[ny'_0]),\qquad g_0(y_0)=nf_0(x_0\otimes y_0).\]

 e）现在使用 $x$ 是 ${}_nP$ 的\Term{section}{截面}这一事实，证明上述 $\operatorname{Ext}^1(nQ,G)$ 的元素是（2.3.15）所定义元素的负元。注意，对任意\Term{homomorphism}{同态} $h:M_0\longrightarrow C^0$，若把 $g_1,g_0$ 分别换成
\[g'_1=g_1+hd_1,\qquad g'_0=g_0-d^0h,\]
则 d）中考虑的 $\operatorname{Ext}^1$ 的元素\Term{invariant}{不变}。而 c）中已构造了\Term{homomorphism}{同态} $\psi:M_0\longrightarrow C^0$；按定义，对 $M'_0$ 的\Term{section}{截面} $y'_0$ 有
\[\psi(y'_0)=f_1([nx_0]\otimes y'_0-x_0\otimes[ny'_0]),\]
也就是
\[f_1(x_0\otimes[ny'_0])\ (=g_1(y'_0))
=f_1([nx_0]\otimes y'_0)-\psi(y'_0).\]
因此定义 $h:M_0\longrightarrow C^0$ 为
\SourcePage{258}
\[h(y_0)=-f_1([nx_0]\otimes y_0)+\psi(y_0).\]
依构造就有
\[g'_1=0,\qquad g'_0(y_0)=nf_0(x_0\otimes y_0)+d^0f_1([nx_0]\otimes y_0)-d^0\psi(y_0).\]
而 $d^0f_1=-f_0d_1$，故 $d^0f_1([nx_0]\otimes y_0)=-f_0(nx_0\otimes y_0)$，从而只剩
\[g'_1=0,\qquad g'_0=-d^0\psi.\]
因此，所考虑的元素由\Term{homomorphism}{同态} $M_0/M'_0\simeq nQ\longrightarrow C^1$ 定义，它是 $-d^0\psi$ 通过取\Term{quotient}{商}得到的。我们得到的正是（2.3.15）所定义元素的负元。
\end{proof}

\textbf{注记 2.3.16.} 上述证明实际上给出了比2.3.11所述略精确的相容性：选取 ${}_nP$ 的一个\Term{section}{截面} $x$，它便定义\Term{homomorphism}{同态} $y\mapsto\varphi_\xi(x,y)$，从 ${}_nQ$ 到 $G$；再经关于（2.3.6）的\Term{coboundary-homomorphism}{余边同态}得到 $\operatorname{Ext}^1(nQ,G)$ 的一个元素。另一方面，暂不考虑 $x$ 被 $n$ \Term{annihilate}{零化}这一事实，取 $x$ 所定义的\Term{homomorphism}{同态} $\mathbf Z_T\longrightarrow P$，通过\Term{inverse-image}{逆像}得到 $(\mathbf Z_T,Q)$ 的一个以 $G$ 为\Term{kernel}{核}的\Term{biextension}{双扩张}；它的类\Term{solvable}{可解}释为 $\operatorname{Ext}^1(Q,G)$ 的一个元素，从而诱导 $\operatorname{Ext}^1(nQ,G)$ 的一个元素。这样得到的两个 $\operatorname{Ext}^1(nQ,G)$ 的元素互为负元。

\textbf{2.3.17.} 利用（2.3.12）中\Term{pairing}{配对} $\varphi_\xi^n:{}_nP\otimes{}_nQ\longrightarrow G$ 的\Term{cocycle}{余圈}描述，再给出此\Term{pairing}{配对}的另一表达式。先假设类为 $\xi$ 的\Term{biextension}{双扩张} $E$ 在 $P\times Q$ 上有\Term{section}{截面} $s$。
\SourcePage{259}
如第VII讲3.8.13，取 $P$ 的\Term{flat-resolution}{平坦解消}为 $L_0=\mathbf Z[P]$，$L_1$ 是\Term{canonical}{典范}\Term{homomorphism}{同态} $\mathbf Z[P]\longrightarrow P$ 的\Term{kernel}{核}；类似定义 $M_\bullet$。那么，$E$ 在 $(L_0,M_0)$ 上的\Term{inverse-image}{逆像}有唯一的\Term{biextension}{双扩张}\Term{trivialization}{平凡化}，使它在 $P\times Q$ 上的\Term{restriction}{限制}由 $s$ 给出；仍记此\Term{trivialization}{平凡化}为 $s$（参见第VII讲3.8.2）。由第VII讲3.8.12，$\xi$ 在 $\operatorname{Ext}^1(P\mathbin{\overset{\mathbf L}{\otimes}}Q,G)$ 中的类由\Term{set-component}{分量}为 $(f,g)$ 的\Term{cocycle}{余圈}给出，其\Term{characteristic}{特征}为
\[s(x_1,y_0)=f(x_1\otimes y_0)e_{E/Q}(y_0),\qquad
s(x_0,y_1)=g(x_0\otimes y_1)e_{E/P}(x_0).\]
若 $x,y$ 分别是 ${}_nP,{}_nQ$ 的\Term{section}{截面}，则将它们\Term{lifting}{提升}为 $L_0,M_0$ 的相应\Term{section}{截面} $x_0=[x],y_0=[y]$，应用（2.3.13）。置 $z=s(x,y)$，得到
\[\varphi_\xi^n(x,y)=\frac{{}^{n}z}{e_{E/Q}(y)}-\frac{z^{n}}{e_{E/P}(x)}.\tag{2.3.18}\]
这里 $z$ 左边（相应地右边）的指数 $n$，表示对 $E$ 上左边（相应地右边）的\Term{multiplication}{乘法}结构取 $n$ 次幂，所得是 $E$ 在 $(0,y)$（相应地 $(x,0)$）上的\Term{section}{截面}；这样的元素除以 $e_{E/Q}(y)$（相应地 $e_{E/P}(x)$）所得“\Term{quotient}{商}”的意义是清楚的。
现在注意，只要 $z$ 是 $E$ 的一个\Term{global-section}{全局截面}，\Term{lifting}{提升}了 ${}_nP,{}_nQ$ 的\Term{section}{截面} $x,y$，公式（2.3.18）仍有意义。我断言，对每个这样的\Term{section}{截面}，（2.3.18）都成立，无须假设 $E$ 在 $P\times Q$ 上有\Term{section}{截面}。因为 $P\times Q$ 的\Term{section}{截面} $(x,y)$ 总能\Term{local-scheme}{局部}\Term{lifting}{提升}为 $E$ 的\Term{section}{截面}（$E\longrightarrow P\times Q$ 是\Term{epimorphism}{满态射}），所以这就在一切情形下给出 $\varphi_\xi^n(x,y)$ 的计算方法。还应注意，若
\SourcePage{260}
将 $z$ 换为 $E$ 的另一个\Term{lifting}{提升} $(x,y)$ 的\Term{section}{截面} $z'$，则 $z'=z+g$，其中 $g$ 是 $G$ 的\Term{section}{截面}；（2.3.18）右边对应于 $z'$ 的两项分别比对应于 $z$ 的项增加 $ng$，所以它们的差确实与 $z$ 无关。这表明（2.3.18）确实定义了一个确定的\Term{homomorphism}{同态} ${}_nP\otimes{}_nQ\longrightarrow G$。
还须证明它就是 $\varphi^n$。为此，立即可归结到 $P=Q=(\mathbf Z/n\mathbf Z)_T$，从而 $P\times Q$ 形如 $I_T$，其中 $I=\mathbf Z/n\mathbf Z\times\mathbf Z/n\mathbf Z$ 是\Term{finite-set}{有限集合}。此时，$G_{P\times Q}$ 下的\Term{torsor}{挠子} $E$（等同于一族 $G$ 下的\Term{torsor}{挠子} $E_i$，$i\in I$）显然在 $T$ 上\Term{locally-trivial}{局部平凡}；因而归结到 $E$ 在 $P\times Q$ 上有\Term{global-section}{全局截面}的情形，这已处理过了。

\textbf{2.3.19. 2.3.11的另一证明（Deligne 给出）。}

 a）归结到 $P=\mathbf Z/n\mathbf Z$，且 $T=(\mathrm{Ens})$。

 b）设 $\partial$ 是如下从 $nQ$ 到 ${}_nQ[1]$ 的\Term{arrow}{箭头}：
\[nQ\xleftarrow{\sim}{[{}_nQ\longrightarrow Q]}\longrightarrow{}_nQ[1].\]
下图\Term{anticommutativity}{反交换}：
\[\begin{tikzcd}
nQ[0]\arrow[r,"1"]\arrow[d,"\partial"']&{[Q\xrightarrow{n}Q]}\\
{}_nQ[1]\arrow[ur,"1"']&
\end{tikzcd}\]
事实上，\Term{arrow}{箭头}
\SourcePage{261}
\[\begin{tikzcd}
{}_nQ\arrow[r]\arrow[d,"0"']&Q\arrow[d,"n"]\\
Q\arrow[r,"n"]&Q
\end{tikzcd}\qquad\text{与}\qquad
\begin{tikzcd}
{}_nQ\arrow[r]\arrow[d,"1"']&Q\arrow[d,"0"]\\
Q\arrow[r,"n"]&Q
\end{tikzcd}\]
互为负元，所用\Term{homotopy}{同伦}为
\[\begin{tikzcd}
{[{}_nQ}\arrow[r]&{Q]}\arrow[dl,"1"']\\
{[Q}\arrow[r,"n"]&{Q]}.
\end{tikzcd}\]

 c）2.3.10就是下图：
\[\resizebox{\linewidth}{!}{\begin{tikzcd}[ampersand replacement=\&,column sep=small,row sep=large]
\&\&H^1\mathbf R\operatorname{Hom}(\mathbf Z\mathbin{\overset{\mathbf L}{\otimes}}Q,G)\arrow[ddr] \&\\
H^1\mathbf R\operatorname{Hom}(\mathbf Z/n\mathbf Z\mathbin{\overset{\mathbf L}{\otimes}}Q,G)\arrow[r]\arrow[d]\arrow[urr]\&
H^1\mathbf R\operatorname{Hom}(\mathbf Z/n\mathbf Z,\mathbf R\mathcal Hom(Q,G))\arrow[d]\&\&\\
\operatorname{Hom}(\mathcal Tor_1(\mathbf Z/n\mathbf Z,Q),G)\arrow[d]\&
\operatorname{Hom}(\mathbf Z/n\mathbf Z,\mathcal Ext^1(Q,G))\arrow[d]\&\&\operatorname{Ext}^1(Q,G)\arrow[dll]\\
\operatorname{Hom}({}_nQ,G)\arrow[r]\&\operatorname{Ext}^1(nQ,G)\&\&
\end{tikzcd}}\]
其中右上部分标记为 $+$，左下部分标记为圈起的 $-\,?$；并且：

 c1）$\mathbf Z/n\mathbf Z\mathbin{\overset{\mathbf L}{\otimes}}Q\simeq[Q\xrightarrow n Q]$；在这个\Term{isomorphism}{同构}下，第一列竖直复合\Term{arrow}{箭头}由
\[{}_nQ[1]\longrightarrow{[Q\xrightarrow n Q]}\]
诱导。

 c2）第二个水平\Term{arrow}{箭头}是上述 b）中的\Term{arrow}{箭头} $\partial$ 的\Term{transpose}{转置}。

 c3）若沿图中上方绕行的路径走，显然得到由 $nQ[0]\longrightarrow[Q\xrightarrow n Q]$ 诱导的\Term{arrow}{箭头}。因此得到2.3.11的结论。

\SourcePage{262}
\textbf{2.4.} 设 $f:T'\longrightarrow T$ 是\Term{topos}{拓扑斯}\Term{morphism}{态射}。则有\Term{canonical}{典范}\Term{homomorphism}{同态}（第VII讲2.8）
\[\operatorname{Biext}^1(P,Q;G)\longrightarrow\operatorname{Biext}^1(P',Q';G'),\]
其中撇号表示\Term{inverse-image}{逆像}\Term{group}{群}。由于对 $T$ 的每个\Term{commutative-group}{交换群} $E$，显然有\Term{canonical-isomorphism}{典范同构}
\[({}_nE)'\simeq{}_n(E'),\]
利用 $T$ 及 $T'$ 上的\Term{homomorphism}{同态}（2.2.2），可定义\Term{canonical}{典范}图
\[\begin{tikzcd}
\operatorname{Biext}^1(P,Q;G)\arrow[r]\arrow[d,"{\varphi^n_{P,Q}}"']&
\operatorname{Biext}^1(P',Q';G')\arrow[d,"{\varphi^n_{P',Q'}}"]\\
\operatorname{Hom}({}_nP\otimes{}_nQ,G)\arrow[r]&\operatorname{Hom}({}_nP'\otimes{}_nQ',G').
\end{tikzcd}\tag{2.4.1}\]
同样，\Term{homomorphism}{同态}（2.2.3）用 $T_\ell$ 给出类似的图。我断言，这些图都是交换的。从定义出发\Term{kernel}{核}实这一点本质上是平凡的，留给读者。更\Term{generic}{一般}地，应当注意，本节及上一节考虑的所有\Term{canonical}{典范}\Term{homomorphism}{同态}，在显然的意义下都与\Term{topos}{拓扑斯}的变换相容。后文将使用这类相容性而不再另外说明。
